\documentclass[11pt]{article}
\usepackage[margin=1in]{geometry}
\usepackage{amsmath,amssymb,amsthm,mathtools}
\usepackage{graphicx}
\usepackage{xcolor}
\usepackage{microtype}
\usepackage[hidelinks]{hyperref}
\usepackage{enumitem}

\newtheorem{theorem}{Theorem}
\newtheorem{lemma}[theorem]{Lemma}
\newcommand{\Bclass}{\mathcal B}
\newcommand{\NCFB}{\mathcal{NCFB}}
\newcommand{\CFB}{\mathcal{CFB}}
\newcommand{\D}{\mathbb D}
\newcommand{\K}{\mathcal K}

\title{A restriction argument for the connected-domain\\
construction question in $\NCFB_n(\Omega)$}
\author{Candidate solution packet for arXiv:1910.10538}
\date{27 August 2026}

\begin{document}
\maketitle

\begin{center}
\fcolorbox{black}{gray!8}{\parbox{0.91\textwidth}{
\textbf{Status.} Candidate full solution to Question 1 as literally stated,
subject to human review.  The construction retains strong irreducibility, but
does not retain the additional spectral equalities in Theorem 3.10 of the
source paper.}}
\end{center}

\begin{abstract}
Jiang, Ji, and Wu ask whether one can construct an operator in
$\NCFB_n(\Omega)$ when the bounded planar domain $\Omega$ is connected rather
than simply connected.  The answer to this literal construction question is
yes.  Take the paper's strongly irreducible example on a disk containing
$\Omega$ and restrict its Cowen--Douglas base domain to $\Omega$.  The only
point needing proof is that the spanning axiom in the Cowen--Douglas
definition survives restriction to a nonempty open subset; this follows from
the identity theorem for a local holomorphic frame.  All additional
$\CFB/\NCFB$ conditions are properties of the same operator and block
decomposition, and hence are unchanged.
\end{abstract}

\section{Source question}

The source is C.~Jiang, K.~Ji, and J.~Wu, \emph{Similarity Invariants of
Essentially normal Cowen--Douglas Operators and Chern Polynomials},
arXiv:1910.10538v1, Section~6, PDF page~27 \cite{JJW}.  The question is:
\begin{quote}
``Can we construct an operator in $\NCFB_n(\Omega)$ when $\Omega$ is
connected?''
\end{quote}
The source context assumes throughout that $\Omega$ is a bounded, open,
connected subset of $\mathbb C$.

\begin{figure}[ht]
  \centering
  \includegraphics[width=\textwidth]{figures/open_problem_crop.png}
  \caption{The two questions at the end of the source paper.  This packet
  answers Question 1.}
\end{figure}

\section{Definitions and proof intuition}

For a connected open set $U\subset\mathbb C$, an operator $A$ belongs to
$\Bclass_m(U)$ when, for every $w\in U$, the operator $A-wI$ is onto,
$\dim\ker(A-wI)=m$, and
\[
 \overline{\operatorname{span}}\{\ker(A-wI):w\in U\}=\mathcal H.
\]
The source defines $\NCFB_n(U)$ as a subclass of the essentially normal
operators in $\CFB_n(U)$, with an upper-triangular block decomposition and
specified idempotents whose Calkin images are self-adjoint \cite{JJW}.

\paragraph{Proof intuition.}
Cowen--Douglas membership is local except for the spanning axiom.  If a vector
is orthogonal to every eigenvector over a nonempty open subset, pair it with a
local holomorphic frame.  The resulting scalar holomorphic functions vanish
on an open set, so the identity theorem propagates their vanishing throughout
the larger connected domain.  Thus the smaller family of eigenspaces is
already total.  Starting with the source paper's disk example, we may
therefore forget all eigenvalues outside $\Omega$ without changing the
operator.  The flag relations, Property~(H), essential normality, and the
idempotents do not depend on which base domain we use.

\section{The restriction lemma}

\begin{lemma}[Restriction of the Cowen--Douglas base]
\label{lem:restriction}
Let $U\subset\mathbb C$ be connected and open, let $V\subset U$ be nonempty
and open, and let $A\in\Bclass_m(U)$.  If $V$ is connected (as in the usual
definition of a Cowen--Douglas domain), then $A\in\Bclass_m(V)$.
\end{lemma}

\begin{proof}
For $w\in V$, surjectivity of $A-wI$ and the equality
$\dim\ker(A-wI)=m$ follow immediately from $A\in\Bclass_m(U)$.  It remains to
prove totality.

Suppose $h\in\mathcal H$ is orthogonal to $\ker(A-wI)$ for every $w\in V$.
Fix $w_0\in V$ and $u\in U$.  Since an open connected subset of $\mathbb C$
is path connected, choose a path in $U$ from $w_0$ to $u$.  Cover its compact
image by a finite chain of open disks
\[
 W_0,W_1,\ldots,W_N\Subset U
\]
such that $W_0\cap V$ is nonempty and open, $u\in W_N$, and every
$W_{j-1}\cap W_j$ is nonempty and open.  Shrinking the disks if necessary,
the Cowen--Douglas bundle has a holomorphic frame
$\gamma_{j,1},\ldots,\gamma_{j,m}$ on each $W_j$.

On $W_0$, each scalar holomorphic function
$w\mapsto\langle\gamma_{0,k}(w),h\rangle$ vanishes on $W_0\cap V$.
The identity theorem makes it vanish on all of $W_0$.  Consequently $h$ is
orthogonal to every fiber over $W_0\cap W_1$.  Applying the identity theorem
to a frame on $W_1$ gives the same conclusion throughout $W_1$.  Induction
along the chain shows that $h\perp\ker(A-uI)$.  As $u\in U$ was arbitrary,
$h$ is orthogonal to all Cowen--Douglas fibers over $U$; their totality gives
$h=0$.  Hence the fibers over $V$ are total.
\end{proof}

\section{Full answer to Question 1}

\begin{theorem}
\label{thm:main}
Let $\Omega\subset\mathbb C$ be a bounded, nonempty, connected open set and
let $n\geq1$.  Then $\NCFB_n(\Omega)$ is nonempty.  In fact, it contains a
strongly irreducible operator.
\end{theorem}

\begin{proof}
The unit-disk part of Theorem~3.10 in the source paper constructs a strongly
irreducible operator $A\in\NCFB_n(\D)$ \cite{JJW}.  Choose $b\in\mathbb C$ and $R>0$ such that
$\Omega\subset U:=b+R\D$, and set
\[
 T=bI+RA.
\]
We first record directly that $T\in\NCFB_n(U)$.  The Cowen--Douglas
conditions transfer under $w=b+Rz$, because
\[
 T-wI=R(A-zI),\qquad \ker(T-wI)=\ker(A-zI).
\]
Essential normality is preserved since
\[
 [T^*,T]=R^2[A^*,A]\in\K(\mathcal H).
\]
Use for $T$ the same upper-triangular block decomposition and idempotents as
for $A$.  If $A_{k,j}=C_{k,j}A_{k,k+1}\cdots A_{j-1,j}$, then for $k<j$
\[
 T_{k,j}=RA_{k,j}
 =R^{1-(j-k)}C_{k,j}T_{k,k+1}\cdots T_{j-1,j},
\]
and the relevant commutants are unchanged by the affine transformation.  For
adjacent diagonal blocks,
\[
 \tau_{T_{j,j},T_{j+1,j+1}}
 =R\,\tau_{A_{j,j},A_{j+1,j+1}},
\]
so both the kernel and the range in Property~(H) are unchanged.  The Calkin
conditions on the idempotents are literally the same.  Thus all defining
$\CFB_n$ and $\NCFB_n$ conditions persist.  Finally,
$\{T\}'=\{A\}'$, so $T$ is strongly irreducible.

Apply Lemma~\ref{lem:restriction} first to $T\in\Bclass_n(U)$ and then to
each diagonal block $T_{j,j}\in\Bclass_1(U)$.  It follows that
$T\in\Bclass_n(\Omega)$ and $T_{j,j}\in\Bclass_1(\Omega)$.  Every remaining
$\CFB_n/\NCFB_n$ condition uses only $T$, its fixed block decomposition,
Property~(H), essential normality, and the same idempotents; none depends on
enlarging or shrinking the Cowen--Douglas base domain.  Therefore
$T\in\NCFB_n(\Omega)$, and its strong irreducibility was already established.
\end{proof}

\section{Verification and limitations}

\begin{itemize}[leftmargin=*]
\item \textbf{Hypothesis check.}  Boundedness of $\Omega$ is part of the
source paper's standing setup and ensures that a containing disk exists.
For an unbounded $\Omega$, no bounded operator can lie in $\Bclass_n(\Omega)$,
since every $w\in\Omega$ would be an eigenvalue.
\item \textbf{Cowen--Douglas check.}  Lemma~\ref{lem:restriction} separately
verifies the only nonlocal axiom, totality.  The other two axioms are inherited
pointwise.
\item \textbf{Structure check.}  The proof explicitly tracks affine scaling
of the upper-triangular products and of the Rosenblum maps used in
Property~(H); the idempotents and their Calkin images do not change.
\item \textbf{What is not claimed.}  The constructed operator has the
spectrum of the containing-disk model, not necessarily
$\overline\Omega$.  Hence this packet does not extend all three extra
conclusions of source Theorem~3.10 to multiply connected domains.  It answers
the final section's literal Question~1, and it does not address Question~2.
\end{itemize}

\section{Bounded novelty check}

On 27 August 2026, the run-wide registry, solution, attempt, and proof-gap
indexes were searched for arXiv:1910.10538, the exact title, $\NCFB_n$, and
connected-domain Cowen--Douglas keywords.  A web/arXiv search used the exact
question and the queries ``NCFB Cowen Douglas connected'', the exact source
title with ``citations'', and ``Cowen--Douglas NCFB connected domain''.  This
located the 2022 journal publication and later papers on Cowen--Douglas flag
classes, but no later source explicitly answering Question~1 or presenting
this restriction argument.  The search is bounded and cannot certify novelty;
the mathematical claim should be reviewed independently, especially because
the argument is short.

\section{Human-review recommendation}

Check three points: (i) the analytic-continuation step in
Lemma~\ref{lem:restriction}; (ii) the scaling of condition~(2) in the
$\CFB_n$ definition and of Property~(H); and (iii) that the intended reading
of Question~1 is its literal nonemptiness/construction statement, rather than
the stronger spectrum-preserving analogue of Theorem~3.10.  If those checks
pass, Theorem~\ref{thm:main} is a complete affirmative answer to Question~1.

\begin{thebibliography}{9}
\bibitem{JJW}
C.~Jiang, K.~Ji, and J.~Wu,
\emph{Similarity Invariants of Essentially normal Cowen--Douglas Operators
and Chern Polynomials}, arXiv:1910.10538v1 (2019); Israel J. Math.
\textbf{248} (2022), 229--270,
\url{https://arxiv.org/abs/1910.10538}.
\end{thebibliography}

\end{document}
